\documentclass[11pt]{article}
\usepackage[margin=1in]{geometry}
\usepackage{amsmath,amssymb}
\usepackage[hidelinks]{hyperref}
\emergencystretch=3em

\newcommand{\Fp}{\mathcal F_p}
\newcommand{\F}{\mathcal F}
\newcommand{\M}{\mathcal M}
\newcommand{\AP}{\mathrm{AP}}

\title{Literature Status for the Discrete Lipschitz-Free Approximation Property Question in arXiv:2005.06555}
\author{Run fa\_banach\_001, lane 6}
\date{June 19, 2026}

\begin{document}
\maketitle

\section*{Status}

This note records that the discrete-metric-space approximation-property
question from Albiac--Ansorena--Cuth--Doucha, arXiv:2005.06555, is answered
positively by Albiac--Ansorena--Bima--Cuth, arXiv:2506.09786.

No new theorem is claimed here. The purpose of this packet is to clear an
open-problem signal from the run queue and identify the exact later theorem
which settles it.

\section*{Source Question}

In the Open Problems section of arXiv:2005.06555, the authors ask:
\begin{quote}
Let $\M$ be a discrete metric space. Does $\F(\M)$ have AP? Or, more
generally, does $\Fp(\M)$ have AP for every $p\in(0,1]$?
\end{quote}
The later paper cites this as AACD21, Question~6.3.

\section*{Later Answer}

The abstract and introduction of arXiv:2506.09786 state that the paper proves
the approximation property for Lipschitz free $p$-spaces over discrete metric
spaces and thereby answers the question raised in arXiv:2005.06555.

The theorem-level statement is arXiv:2506.09786, Theorem~\texttt{thm:discreteCase}:
if $\M$ is a complete metric space with at most finitely many accumulation
points, then $\Fp(\M)$ has the approximation property for every $p\in(0,1]$.
The line immediately before the theorem says that this answers
AACD21, Question~6.3, in the positive.

In the proof of the theorem, the authors invoke Proposition~4.1 of
arXiv:2005.06555, and the comment before Question~6.3 there, to reduce to the
case where $\M$ is a complete discrete metric space. They then use their
compact-reduction principle and finite-rank approximants to prove the
approximation property. Thus the source question is not merely implied by an
agent-side reformulation: the supporting paper explicitly labels and answers
the same question.

\section*{Scope}

The resolved scope is exactly the discrete metric space AP question for
$0<p\leq 1$. This packet does not settle the neighboring countable proper
metric approximation property question, the problem asking for a metric space
$\M$ with $\Fp(\M)$ failing AP, or the other isomorphism/basis/net questions
listed in the same Open Problems section of arXiv:2005.06555.

\section*{References}

\begin{enumerate}
\item F. Albiac, J. L. Ansorena, M. Cuth, and M. Doucha,
\emph{Lipschitz free spaces isomorphic to their infinite sums and geometric
applications}, arXiv:2005.06555; Trans. Amer. Math. Soc. 374 (2021),
7281--7312.
\item F. Albiac, J. L. Ansorena, J. Bima, and M. Cuth,
\emph{Lipschitz free $p$-spaces for $0<p<1$ in the light of the Schur
$p$-property and the compact reduction}, arXiv:2506.09786.
\end{enumerate}

\end{document}

